二分查找深度题卡 按题型拆成章内大节，但每个大节都先从 LeetCode 案例切入，再进入分流、案例矩阵、案例推导和实验复盘。这样读者看到的是从具体题推到规律，而不是先读抽象方法再猜它能用在哪。

\topic{二分查找与答案空间}
这一节先看 LeetCode 4 寻找两个正序数组的中位数、LeetCode 33 搜索旋转排序数组、LeetCode 34 排序数组中查找首尾位置 这几道 LeetCode 题，再整理同一题型的分流和推导。阅读顺序不是先背原理，而是先看案例的暴力瓶颈，再把重复出现的观察抽象成方法。

\subsection{原理和适用条件}

以 LeetCode 4 寻找两个正序数组的中位数、LeetCode 33 搜索旋转排序数组、LeetCode 34 排序数组中查找首尾位置 为主线：LeetCode 4 寻找两个正序数组的中位数：模型是“两个有序数组的中位数可以看成左右两半切分后的边界值。”，优化抓手是“二分较短数组的切分位置，让左半元素数量正确且左半最大值不大于右半最小值。”；LeetCode 33 搜索旋转排序数组：模型是“旋转数组每次二分至少有一半保持有序。”，优化抓手是“判断哪一半有序，再看目标是否落在有序半边范围内。”；LeetCode 34 排序数组中查找首尾位置：模型是“首尾位置可以拆成两个边界：第一个大于等于和第一个大于。”，优化抓手是“不要找到一个目标后向两边线性扩展，否则重复元素很多时退化。”。这些案例共同推出的规律是：先线性扫描下标或线性试答案，再找出能排除一侧候选的规则。这个规则可能是有序边界、答案可行性，也可能是旋转数组或峰值里的局部比较；二分的核心不是 mid 写法，而是被排除的一侧确实不可能破坏答案。 方法名只是复盘后的标签，真正要掌握的是从题面约束、暴力失败、状态设计到边界反例的推导链。

\begin{principlebox}{图示：从 LeetCode 案例到 二分查找与答案空间推导链}
\begin{tabular}{@{}p{0.18\linewidth}p{0.74\linewidth}@{}}
暴力基线 & 枚举候选、完整模拟或逐个检查，先保证覆盖全部可能答案。\\
瓶颈定位 & 慢点在于每次只排除一个候选。二分能成立，是因为一次比较可以排除一侧下标，或者一次检查可以排除一段答案值域。\\
结构选择 & 二分前必须先写出可排除规则：有序数组用边界谓词，答案空间用可行性谓词，旋转数组和峰值用局部比较证明某一侧必无答案或必有答案。\\
不变量 & 建议固定左闭右开或左闭右闭中的一种。循环中始终维护答案仍在候选区间内；每次移动 \code{left} 或 \code{right} 前，都要证明被丢弃的一侧不可能含有目标。\\
代码落地 & 中点写成 \code{left + (right - left) / 2}。判断平方、乘法、总量时避免溢出。答案二分把检查函数单独写出；局部比较二分把每个分支的排除理由写清楚。\\
\end{tabular}
\end{principlebox}

\subsection{LeetCode 案例分流}

\begin{principlebox}{从 LeetCode 案例分流：二分查找与答案空间}
\textbf{基础边界查找。}
代表案例：LeetCode 704 二分查找、LeetCode 35 搜索插入位置、LeetCode 34 查找首尾位置。先看这些题的题面条件：数组有序，答案是第一个满足某条件的位置。由此推出的处理方式是：把查找目标改写成 lower bound 或 upper bound。风险点：空数组、重复值和返回边界是关键。

\textbf{旋转和局部有序。}
代表案例：LeetCode 33 搜索旋转排序数组、LeetCode 153 寻找旋转排序数组最小值、LeetCode 162 寻找峰值。先看这些题的题面条件：整体不完全有序，但每次比较能确定一半有序或一侧含答案。由此推出的处理方式是：判断有序半边或趋势方向，再排除不含答案的一侧。风险点：重复元素会破坏一半有序判断，可能退化。

\textbf{答案空间二分。}
代表案例：LeetCode 875 爱吃香蕉、LeetCode 1011 船运包裹、LeetCode 410 分割数组最大值、LeetCode 1631 最小体力路径。先看这些题的题面条件：答案值越大越容易或越难满足要求，存在可行边界。由此推出的处理方式是：定义检查函数，在答案范围上找最小可行或最大可行值。风险点：检查函数本身常包含贪心或搜索，必须单独证明。

\textbf{矩阵和二维排除。}
代表案例：LeetCode 240 搜索二维矩阵 II、LeetCode 74 搜索二维矩阵。先看这些题的题面条件：行列有序或某个角点具备单调排除能力。由此推出的处理方式是：从可排除一行或一列的位置开始扫描，或把矩阵视作有序序列二分。风险点：从错误角落出发无法排除足够候选。

\end{principlebox}

\subsection{案例矩阵}

本节案例覆盖 LeetCode 4 寻找两个正序数组的中位数、LeetCode 33 搜索旋转排序数组、LeetCode 34 排序数组中查找首尾位置、LeetCode 35 搜索插入位置、LeetCode 69 x 的平方根、LeetCode 74 搜索二维矩阵、LeetCode 81 搜索旋转排序数组 II、LeetCode 153 寻找旋转排序数组中的最小值、LeetCode 154 寻找旋转排序数组中的最小值 II、LeetCode 162 寻找峰值、LeetCode 240 搜索二维矩阵 II、LeetCode 315 计算右侧小于当前元素的个数、LeetCode 410 分割数组的最大值、LeetCode 540 有序数组中的单一元素、LeetCode 704 二分查找、LeetCode 875 爱吃香蕉的珂珂、LeetCode 981 基于时间的键值存储、LeetCode 1011 在 D 天内送达包裹、LeetCode 1044 最长重复子串、LeetCode 1482 制作 m 束花所需的最少天数。阅读时不要按题号背答案，而要先判断题面落在哪个分流场景：查询目标是什么，输入约束给了什么单调性或等价关系，暴力慢在哪里，优化结构保存了什么状态。

\begin{principlebox}{案例矩阵}
\textbf{LeetCode 4 寻找两个正序数组的中位数。}
两个有序数组的中位数可以看成左右两半切分后的边界值。单点风险：某个数组为空、总长度奇偶、切分在边界、哨兵值不能参与真实比较。

\textbf{LeetCode 33 搜索旋转排序数组。}
旋转数组每次二分至少有一半保持有序。单点风险：未旋转、长度一、目标不存在、边界等号、存在重复元素时退化。

\textbf{LeetCode 34 排序数组中查找首尾位置。}
首尾位置可以拆成两个边界：第一个大于等于和第一个大于。单点风险：目标不存在、目标在开头或结尾、数组为空、全是目标。

\textbf{LeetCode 35 搜索插入位置。}
题目等价于寻找第一个大于等于目标的位置。单点风险：目标小于所有元素、目标大于所有元素、重复值、空数组。

\textbf{LeetCode 69 x 的平方根。}
答案是最大的整数 r，使得 r 的平方不超过 x。单点风险：x 为零、一、非完全平方、接近 int 最大值。

\textbf{LeetCode 74 搜索二维矩阵。}
矩阵每行有序且下一行首元素大于上一行末元素，可视为一维有序数组。单点风险：空矩阵、单行、单列、目标小于全部或大于全部。

\textbf{LeetCode 81 搜索旋转排序数组 II。}
旋转数组存在重复值时，二分可能无法判断哪一半有序。单点风险：全相同元素、目标不存在、重复值跨过旋转点、退化到线性。

\textbf{LeetCode 153 寻找旋转排序数组中的最小值。}
最小值是旋转断点，右端点可帮助判断中点在哪一段。单点风险：未旋转、长度一、旋转一位、无重复条件。

\textbf{LeetCode 154 寻找旋转排序数组中的最小值 II。}
重复元素会让 mid 和 right 相等时无法判断最小值方向。单点风险：全相等、重复值跨断点、最小值出现多次。

\textbf{LeetCode 162 寻找峰值。}
比较 mid 和 mid 加一 的斜率能判断某侧一定存在峰值。单点风险：长度一、峰在边界、多个峰、相邻不相等假设。

\textbf{LeetCode 240 搜索二维矩阵 II。}
右上角同时具有向左变小、向下变大的排除能力。单点风险：空矩阵、单行、单列、目标在角落、重复值。

\textbf{LeetCode 315 计算右侧小于当前元素的个数。}
从右向左扫描时，已处理元素正好是当前位置右侧元素集合。单点风险：重复值、负数、值域很大、严格小于不能包含相等值、结果顺序要反向写回。

\textbf{LeetCode 410 分割数组的最大值。}
连续分成 m 段后最大段和越大，越容易完成分割。单点风险：m 为一、m 等于数组长度、单个元素超过猜测上限、总和溢出。

\textbf{LeetCode 540 有序数组中的单一元素。}
除一个元素外其余元素都成对出现，单一元素会改变配对下标的奇偶规律。单点风险：单一元素在开头或结尾、数组长度一、mid 加一越界。

\textbf{LeetCode 704 二分查找。}
基础题的重点是固定区间语义。单点风险：空数组、长度一、目标在边界、目标不存在。

\textbf{LeetCode 875 爱吃香蕉的珂珂。}
速度越大，总耗时越小，答案具有单调可行性。单点风险：速度下界、堆很大、h 等于堆数、乘法加法溢出。

\textbf{LeetCode 981 基于时间的键值存储。}
同一个 key 下的记录按时间戳递增，查询是不超过目标时间的最新值。单点风险：key 不存在、查询早于第一条记录、相同 key 多次 set、时间戳递增假设。

\textbf{LeetCode 1011 在 D 天内送达包裹。}
容量越大，需要天数越少，存在最小可行容量。单点风险：D 为一、D 等于包裹数、单个包裹超过猜测容量。

\textbf{LeetCode 1044 最长重复子串。}
重复子串长度具有可行性单调：存在长度 L 的重复，则更短长度也存在。单点风险：全相同字符、无重复、哈希冲突、长度边界、模数选择。

\textbf{LeetCode 1482 制作 m 束花所需的最少天数。}
天数越大，可用花越多，能制作的相邻花束数不会减少。单点风险：m*k 超过花数、k 为一、所有花同日开放、相邻性、乘法溢出。

\end{principlebox}

\subsection{共性落点}

\begin{principlebox}{本组共性落点}
\textbf{慢版本。} 暴力做法通常线性扫描下标、逐个试答案值，或直接检查所有候选。它的问题不是不会找到答案，而是没有利用输入顺序、局部比较或答案可行性的边界信息。
\textbf{瓶颈。} 慢点在于每次只排除一个候选。二分能成立，是因为一次比较可以排除一侧下标，或者一次检查可以排除一段答案值域。
\textbf{状态字段。} left、right、mid、mid 处比较值或检查返回值、答案区间含义、返回值语义；每一轮更新都要保留目标位置或第一个可行值。
\textbf{不变量。} 建议固定左闭右开或左闭右闭中的一种。循环中始终维护答案仍在候选区间内；每次移动 \code{left} 或 \code{right} 前，都要证明被丢弃的一侧不可能含有目标。
\textbf{实现检查。} 检查函数或比较分支单独写清，mid 不溢出，循环退出条件和返回值配套；每个分支都要解释丢弃区间。
\textbf{C++ 落地。} 中点写成 \code{left + (right - left) / 2}。判断平方、乘法、总量时避免溢出。答案二分把检查函数单独写出；局部比较二分把每个分支的排除理由写清楚。
\textbf{证明抓手。} 证明从排除一侧开始：答案空间题证明谓词单调，局部比较题证明保留下来的区间仍含目标；不能只靠经验猜测左右边界。
\end{principlebox}

\subsection{案例推导}

\noindent\textbf{LeetCode 4 寻找两个正序数组的中位数。}

\textbf{题目说明。} LeetCode 4 寻找两个正序数组的中位数 的题面入口要先还原成具体任务：两个有序数组的中位数可以看成左右两半切分后的边界值。

\textbf{样例入口。} 拿 [1, 3] 和 [2] 手算，合并后中位数是 2；但不用真合并，只要把总长度切成左半两个数、右半一个数。再拿 [1, 2] 和 [3, 4]，如果短数组切 1 个、长数组切 1 个，左半最大 3 大于右半最小 2，切分错了，应把短数组切分往右调。

\textbf{暴力基线。} 慢版本按顺序扫描下标、逐个试答案值，或直接检查候选直到遇到目标边界。它先保证候选空间覆盖完整，但每次只排除一个候选，浪费了题目给出的顺序、比较或可行性边界。放到本题就是：先按“两个有序数组的中位数可以看成左右两半切分后的边界值”完整试慢版本；样例里已经能看到“规律是二分短数组切分点，使左半元素个数正确且左右边界有序”，所以优化前必须先能说清慢版本枚举了哪些候选。

\textbf{暴力过程。} 拿 [1, 3] 和 [2] 手算，合并后中位数是 2；但不用真合并，只要把总长度切成左半两个数、右半一个数。再拿 [1, 2] 和 [3, 4]，如果短数组切 1 个、长数组切 1 个，左半最大 3 大于右半最小 2，切分错了，应把短数组切分往右调。

\textbf{重复工作。} 题面模型：两个有序数组的中位数可以看成左右两半切分后的边界值。暴力版的慢点要落到本题：慢版本会按顺序试“两个有序数组的中位数可以看成左右两半切分后的边界值”里的候选；而样例规律已经指出“规律是二分短数组切分点，使左半元素个数正确且左右边界有序”。一次比较或一次可行性检查能丢掉一整段候选，继续线性试就是浪费。后面的优化只消掉这类重复工作，不能改变答案集合。

\textbf{优化条件。} 本题的关键观察是：二分较短数组的切分位置，让左半元素数量正确且左半最大值不大于右半最小值。这句话必须能翻译成可维护状态；如果题目条件改变后观察不再成立，就不能继续套原来的写法。

\textbf{相邻状态。} 规律是二分短数组切分点，使左半元素个数正确且左右边界有序。把这个特例继续拆开看：每次比较 mid 之后，必须能指出哪一侧整体不可能包含目标，或哪一侧整体已经满足可行性。二分的核心不是 mid 公式，而是被丢弃区间确实不含答案。

\textbf{状态复用。} 把逐个试候选改成维护答案所在区间；边界谓词、可行性检查或局部比较给出分支方向，每个分支都必须能排除一侧候选。本题落点是：规律是二分短数组切分点，使左半元素个数正确且左右边界有序。手算时只改被新输入影响的那一小部分，而不是回头重算完整候选。

\textbf{指针和字段。} 状态字段要能说清：left、right、mid、mid 处比较值或检查返回值、答案区间含义、返回值语义；每一轮更新都要保留目标位置或第一个可行值。手算时按这个协议跟踪：拿 [1, 3] 和 [2] 手算，合并后中位数是 2；但不用真合并，只要把总长度切成左半两个数、右半一个数。再拿 [1, 2] 和 [3, 4]，如果短数组切 1 个、长数组切 1 个，左半最大 3 大于右半最小 2，切分错了，应把短数组切分往右调。然后按协议复查：手算时列出 left、mid、right，以及 mid 处比较结果、可行性检查结果或局部趋势。每一步都要写出被丢弃的一侧为什么不含答案。

\textbf{代码落点。} 检查函数或比较分支单独写清，mid 不溢出，循环退出条件和返回值配套；每个分支都要解释丢弃区间。关键代码行必须能逐句翻译成“旧状态怎样变成新状态”，否则就只是模板。

\textbf{正确性。} 证明从排除一侧开始：答案空间题证明谓词单调，局部比较题证明保留下来的区间仍含目标；不能只靠经验猜测左右边界。

\textbf{边界实验。} 先用某个数组为空、总长度奇偶、切分在边界、哨兵值不能参与真实比较做反例检查。实验时覆盖某个数组为空、总长度奇偶、切分在边界、哨兵值不能参与真实比较，逐轮打印 left、mid、right、比较值或检查结果。只要有一次被排除区间仍含答案，二分不变量就错了。

\textbf{复杂度变化。} 暴力逐个试候选是线性或和值域成正比；二分每次排除一半候选，复杂度变成对数级，答案二分还要乘上单次检查成本。

\textbf{迁移追问。} 如果题目改成“若只要求第 k 小元素，可以用递归丢弃较小前缀或二分切分的同类思想”，先判断原状态是否仍然覆盖新答案集合，再决定是微调边界、改 value 语义，还是换算法。

\noindent\textbf{LeetCode 33 搜索旋转排序数组。}

\textbf{题目说明。} LeetCode 33 搜索旋转排序数组 的题面入口要先还原成具体任务：旋转数组每次二分至少有一半保持有序。

\textbf{样例入口。} 拿 [4, 5, 6, 7, 0, 1, 2] 找 0 手算，mid=7 时左半 [4, 5, 6, 7] 有序，但目标 0 不在这个范围，所以整段左半都能丢掉。下一轮目标落在右半。

\textbf{暴力基线。} 慢版本按顺序扫描下标、逐个试答案值，或直接检查候选直到遇到目标边界。它先保证候选空间覆盖完整，但每次只排除一个候选，浪费了题目给出的顺序、比较或可行性边界。放到本题就是：先按“旋转数组每次二分至少有一半保持有序”完整试慢版本；样例里已经能看到“规律是每次先判断哪一半有序，再用目标是否落入有序半边来排除一整段；有重复元素时有序半边可能判断不出来”，所以优化前必须先能说清慢版本枚举了哪些候选。

\textbf{暴力过程。} 拿 [4, 5, 6, 7, 0, 1, 2] 找 0 手算，mid=7 时左半 [4, 5, 6, 7] 有序，但目标 0 不在这个范围，所以整段左半都能丢掉。下一轮目标落在右半。

\textbf{重复工作。} 题面模型：旋转数组每次二分至少有一半保持有序。暴力版的慢点要落到本题：慢版本会按顺序试“旋转数组每次二分至少有一半保持有序”里的候选；而样例规律已经指出“规律是每次先判断哪一半有序，再用目标是否落入有序半边来排除一整段；有重复元素时有序半边可能判断不出来”。一次比较或一次可行性检查能丢掉一整段候选，继续线性试就是浪费。后面的优化只消掉这类重复工作，不能改变答案集合。

\textbf{优化条件。} 本题的关键观察是：判断哪一半有序，再看目标是否落在有序半边范围内。这句话必须能翻译成可维护状态；如果题目条件改变后观察不再成立，就不能继续套原来的写法。

\textbf{相邻状态。} 规律是每次先判断哪一半有序，再用目标是否落入有序半边来排除一整段；有重复元素时有序半边可能判断不出来。把这个特例继续拆开看：每次比较 mid 之后，必须能指出哪一侧整体不可能包含目标，或哪一侧整体已经满足可行性。二分的核心不是 mid 公式，而是被丢弃区间确实不含答案。

\textbf{状态复用。} 把逐个试候选改成维护答案所在区间；边界谓词、可行性检查或局部比较给出分支方向，每个分支都必须能排除一侧候选。本题落点是：规律是每次先判断哪一半有序，再用目标是否落入有序半边来排除一整段；有重复元素时有序半边可能判断不出来。手算时只改被新输入影响的那一小部分，而不是回头重算完整候选。

\textbf{指针和字段。} 状态字段要能说清：left、right、mid、mid 处比较值或检查返回值、答案区间含义、返回值语义；每一轮更新都要保留目标位置或第一个可行值。手算时按这个协议跟踪：拿 [4, 5, 6, 7, 0, 1, 2] 找 0 手算，mid=7 时左半 [4, 5, 6, 7] 有序，但目标 0 不在这个范围，所以整段左半都能丢掉。下一轮目标落在右半。然后按协议复查：手算时列出 left、mid、right，以及 mid 处比较结果、可行性检查结果或局部趋势。每一步都要写出被丢弃的一侧为什么不含答案。

\textbf{代码落点。} 检查函数或比较分支单独写清，mid 不溢出，循环退出条件和返回值配套；每个分支都要解释丢弃区间。关键代码行必须能逐句翻译成“旧状态怎样变成新状态”，否则就只是模板。

\textbf{正确性。} 证明从排除一侧开始：答案空间题证明谓词单调，局部比较题证明保留下来的区间仍含目标；不能只靠经验猜测左右边界。

\textbf{边界实验。} 先用未旋转、长度一、目标不存在、边界等号、存在重复元素时退化做反例检查。实验时覆盖未旋转、长度一、目标不存在、边界等号、存在重复元素时退化，逐轮打印 left、mid、right、比较值或检查结果。只要有一次被排除区间仍含答案，二分不变量就错了。

\textbf{复杂度变化。} 暴力逐个试候选是线性或和值域成正比；二分每次排除一半候选，复杂度变成对数级，答案二分还要乘上单次检查成本。

\textbf{迁移追问。} 如果题目改成“若重复元素很多，需要在无法判断有序半边时收缩边界”，先判断原状态是否仍然覆盖新答案集合，再决定是微调边界、改 value 语义，还是换算法。

\noindent\textbf{LeetCode 34 排序数组中查找首尾位置。}

\textbf{题目说明。} LeetCode 34 排序数组中查找首尾位置 的题面入口要先还原成具体任务：首尾位置可以拆成两个边界：第一个大于等于和第一个大于。

\textbf{样例入口。} 拿 [5, 7, 7, 8, 8, 10] 找 8 手算，找到一个 8 后向两边扩展能得到答案，但全数组都是 8 时会退化成线性。改成两次边界二分：第一个 >=8 的位置是 3，第一个 >8 的位置是 5，区间就是 [3, 4]。

\textbf{暴力基线。} 暴力线性扫描数组，记录第一个等于 target 的下标和最后一个等于 target 的下标。

\textbf{暴力过程。} 以 nums=[5, 7, 7, 8, 8, 10]、target=8 为例，线性扫描会在下标 3 第一次看到 8，在下标 4 最后一次看到 8。

\textbf{重复工作。} 题面模型：首尾位置可以拆成两个边界：第一个大于等于和第一个大于。暴力版的慢点要落到本题：浪费在数组已经有序却仍然逐个看。目标值出现的位置一定是一段连续区间，只需要找这段区间的左右边界。后面的优化只消掉这类重复工作，不能改变答案集合。

\textbf{优化条件。} 本题的关键观察是：不要找到一个目标后向两边线性扩展，否则重复元素很多时退化。这句话必须能翻译成可维护状态；如果题目条件改变后观察不再成立，就不能继续套原来的写法。

\textbf{相邻状态。} 规律是首尾位置不是一次查找，而是两个单调边界。把这个特例继续拆开看：每次比较 mid 之后，必须能指出哪一侧整体不可能包含目标，或哪一侧整体已经满足可行性。二分的核心不是 mid 公式，而是被丢弃区间确实不含答案。

\textbf{状态复用。} 把问题拆成两个 lower\_bound：第一个大于等于 target 的位置 left，第一个大于 target 的位置 right\_exclusive。手算时只改被新输入影响的那一小部分，而不是回头重算完整候选。

\textbf{指针和字段。} 状态字段要能说清：left\_bound(x) 返回第一个 nums[i]>=x 的下标；答案是 left=left\_bound(target)，right=left\_bound(target+1)-1。手算时按这个协议跟踪：样例中 lower\_bound(8)=3，lower\_bound(9)=5，所以答案是 [3, 4]；若 lower\_bound(target) 越界或值不等于 target，返回 [-1, -1]。

\textbf{代码落点。} 关键是边界函数使用左闭右开 [left, right)，循环条件 left<right；退出时 left 就是第一个满足谓词的位置。关键代码行必须能逐句翻译成“旧状态怎样变成新状态”，否则就只是模板。

\textbf{正确性。} 有序数组中 nums[i]>=x 是单调谓词，二分能找到第一个 true；target 的所有位置正好夹在第一个 >=target 和第一个 >target 之间。

\textbf{边界实验。} 先用目标不存在、目标在开头或结尾、数组为空、全是目标做反例检查。实验时覆盖目标不存在、目标在开头或结尾、数组为空、全是目标，逐轮打印 left、mid、right、比较值或检查结果。只要有一次被排除区间仍含答案，二分不变量就错了。

\textbf{复杂度变化。} 线性扫描 O(n)，两次二分 O(log n)，空间 O(1)。

\textbf{迁移追问。} 如果题目改成“这个模型能迁移到计数某值出现次数和插入区间边界”，先判断原状态是否仍然覆盖新答案集合，再决定是微调边界、改 value 语义，还是换算法。

\noindent\textbf{LeetCode 35 搜索插入位置。}

\textbf{题目说明。} LeetCode 35 搜索插入位置 的题面入口要先还原成具体任务：题目等价于寻找第一个大于等于目标的位置。

\textbf{样例入口。} 拿 [1, 3, 5, 6] 插入 2 手算，线性看 1 还小、3 已经不小，所以答案是下标 1。二分时保持 left 是第一个可能不小于目标的位置，mid 值小于目标才丢左侧。

\textbf{暴力基线。} 暴力从左到右扫描，找到第一个 nums[i]>=target 的位置；若一直找不到，插入位置就是 n。

\textbf{暴力过程。} 以 nums=[1, 3, 5, 6]、target=2 为例，1 小于 2，3 已经不小于 2，所以答案是 1；target=7 时扫描到末尾后返回 4。

\textbf{重复工作。} 题面模型：题目等价于寻找第一个大于等于目标的位置。暴力版的慢点要落到本题：浪费在有序数组中一次只排除一个元素。nums[mid] 小于 target 时，mid 及其左侧都不可能是插入位置。后面的优化只消掉这类重复工作，不能改变答案集合。

\textbf{优化条件。} 本题的关键观察是：统一用 lower bound 语义，不在相等时提前返回也能保持边界一致。这句话必须能翻译成可维护状态；如果题目条件改变后观察不再成立，就不能继续套原来的写法。

\textbf{相邻状态。} 规律是统一成 lower\_bound，不需要在相等时提前返回；目标大于所有元素时 left 会走到 n。把这个特例继续拆开看：每次比较 mid 之后，必须能指出哪一侧整体不可能包含目标，或哪一侧整体已经满足可行性。二分的核心不是 mid 公式，而是被丢弃区间确实不含答案。

\textbf{状态复用。} 二分第一个 nums[i]>=target 的下标，也就是 lower\_bound。候选区间始终包含插入位置。手算时只改被新输入影响的那一小部分，而不是回头重算完整候选。

\textbf{指针和字段。} 状态字段要能说清：left 是当前可能答案的左边界，right 是开区间右边界，mid 用来判断左半是否全部小于 target。手算时按这个协议跟踪：[1, 3, 5, 6] 找 2：mid 指向 5 时右边界左移；mid 指向 3 时继续左移；mid 指向 1 时左边界右移，最终 left=1。

\textbf{代码落点。} 关键分支是 nums[mid] < target 时 left=mid+1，否则 right=mid。不要把等于 target 提前返回成特殊分支，统一边界更稳。关键代码行必须能逐句翻译成“旧状态怎样变成新状态”，否则就只是模板。

\textbf{正确性。} 被 left 跳过的元素都严格小于 target，不可能成为第一个不小于 target 的位置；right 左侧收缩时保留了可能答案。

\textbf{边界实验。} 先用目标小于所有元素、目标大于所有元素、重复值、空数组做反例检查。实验时覆盖目标小于所有元素、目标大于所有元素、重复值、空数组，逐轮打印 left、mid、right、比较值或检查结果。只要有一次被排除区间仍含答案，二分不变量就错了。

\textbf{复杂度变化。} 暴力 O(n)，二分 O(log n)，空间 O(1)。

\textbf{迁移追问。} 如果题目改成“手写二分时用左闭右开可直接返回 left”，先判断原状态是否仍然覆盖新答案集合，再决定是微调边界、改 value 语义，还是换算法。

\noindent\textbf{LeetCode 69 x 的平方根。}

\textbf{题目说明。} LeetCode 69 x 的平方根 的题面入口要先还原成具体任务：答案是最大的整数 r，使得 r 的平方不超过 x。

\textbf{样例入口。} 拿 x=8 手算，二的平方不超过 8，而三的平方已经超过 8，所以答案是 2。

\textbf{暴力基线。} 慢版本按顺序扫描下标、逐个试答案值，或直接检查候选直到遇到目标边界。它先保证候选空间覆盖完整，但每次只排除一个候选，浪费了题目给出的顺序、比较或可行性边界。放到本题就是：先按“答案是最大的整数 r，使得 r 的平方不超过 x”完整试慢版本；样例里已经能看到“规律是找最后一个平方不超过 x 的整数，比较平方时用 long long 或除法避免溢出”，所以优化前必须先能说清慢版本枚举了哪些候选。

\textbf{暴力过程。} 拿 x=8 手算，二的平方不超过 8，而三的平方已经超过 8，所以答案是 2。

\textbf{重复工作。} 题面模型：答案是最大的整数 r，使得 r 的平方不超过 x。暴力版的慢点要落到本题：慢版本会按顺序试“答案是最大的整数 r，使得 r 的平方不超过 x”里的候选；而样例规律已经指出“规律是找最后一个平方不超过 x 的整数，比较平方时用 long long 或除法避免溢出”。一次比较或一次可行性检查能丢掉一整段候选，继续线性试就是浪费。后面的优化只消掉这类重复工作，不能改变答案集合。

\textbf{优化条件。} 本题的关键观察是：在整数范围上二分右边界，比较时用除法或 long long 避免平方溢出。这句话必须能翻译成可维护状态；如果题目条件改变后观察不再成立，就不能继续套原来的写法。

\textbf{相邻状态。} 规律是找最后一个平方不超过 x 的整数，比较平方时用 long long 或除法避免溢出。把这个特例继续拆开看：每次比较 mid 之后，必须能指出哪一侧整体不可能包含目标，或哪一侧整体已经满足可行性。二分的核心不是 mid 公式，而是被丢弃区间确实不含答案。

\textbf{状态复用。} 把逐个试候选改成维护答案所在区间；边界谓词、可行性检查或局部比较给出分支方向，每个分支都必须能排除一侧候选。本题落点是：规律是找最后一个平方不超过 x 的整数，比较平方时用 long long 或除法避免溢出。手算时只改被新输入影响的那一小部分，而不是回头重算完整候选。

\textbf{指针和字段。} 状态字段要能说清：left、right、mid、mid 处比较值或检查返回值、答案区间含义、返回值语义；每一轮更新都要保留目标位置或第一个可行值。手算时按这个协议跟踪：拿 x=8 手算，二的平方不超过 8，而三的平方已经超过 8，所以答案是 2。然后按协议复查：手算时列出 left、mid、right，以及 mid 处比较结果、可行性检查结果或局部趋势。每一步都要写出被丢弃的一侧为什么不含答案。

\textbf{代码落点。} 检查函数或比较分支单独写清，mid 不溢出，循环退出条件和返回值配套；每个分支都要解释丢弃区间。关键代码行必须能逐句翻译成“旧状态怎样变成新状态”，否则就只是模板。

\textbf{正确性。} 证明从排除一侧开始：答案空间题证明谓词单调，局部比较题证明保留下来的区间仍含目标；不能只靠经验猜测左右边界。

\textbf{边界实验。} 先用x 为零、一、非完全平方、接近 int 最大值做反例检查。实验时覆盖x 为零、一、非完全平方、接近 int 最大值，逐轮打印 left、mid、right、比较值或检查结果。只要有一次被排除区间仍含答案，二分不变量就错了。

\textbf{复杂度变化。} 暴力逐个试候选是线性或和值域成正比；二分每次排除一半候选，复杂度变成对数级，答案二分还要乘上单次检查成本。

\textbf{迁移追问。} 如果题目改成“若要求浮点精度，可以二分实数区间或用牛顿迭代”，先判断原状态是否仍然覆盖新答案集合，再决定是微调边界、改 value 语义，还是换算法。

\noindent\textbf{LeetCode 74 搜索二维矩阵。}

\textbf{题目说明。} LeetCode 74 搜索二维矩阵 的题面入口要先还原成具体任务：矩阵每行有序且下一行首元素大于上一行末元素，可视为一维有序数组。

\textbf{样例入口。} 拿矩阵 [[1, 3, 5], [7, 9, 11]] 找 9 手算，把它看成一维有序数组 [1, 3, 5, 7, 9, 11]，下标 4 映射到 row=1、col=1。

\textbf{暴力基线。} 慢版本按顺序扫描下标、逐个试答案值，或直接检查候选直到遇到目标边界。它先保证候选空间覆盖完整，但每次只排除一个候选，浪费了题目给出的顺序、比较或可行性边界。放到本题就是：先按“矩阵每行有序且下一行首元素大于上一行末元素，可视为一维有序数组”完整试慢版本；样例里已经能看到“规律是行尾小于下一行行首时可以整体展平二分，空矩阵和单行要先处理”，所以优化前必须先能说清慢版本枚举了哪些候选。

\textbf{暴力过程。} 拿矩阵 [[1, 3, 5], [7, 9, 11]] 找 9 手算，把它看成一维有序数组 [1, 3, 5, 7, 9, 11]，下标 4 映射到 row=1、col=1。

\textbf{重复工作。} 题面模型：矩阵每行有序且下一行首元素大于上一行末元素，可视为一维有序数组。暴力版的慢点要落到本题：慢版本会按顺序试“矩阵每行有序且下一行首元素大于上一行末元素，可视为一维有序数组”里的候选；而样例规律已经指出“规律是行尾小于下一行行首时可以整体展平二分，空矩阵和单行要先处理”。一次比较或一次可行性检查能丢掉一整段候选，继续线性试就是浪费。后面的优化只消掉这类重复工作，不能改变答案集合。

\textbf{优化条件。} 本题的关键观察是：把下标 mid 映射到 row 和 col，再做标准二分。这句话必须能翻译成可维护状态；如果题目条件改变后观察不再成立，就不能继续套原来的写法。

\textbf{相邻状态。} 规律是行尾小于下一行行首时可以整体展平二分，空矩阵和单行要先处理。把这个特例继续拆开看：每次比较 mid 之后，必须能指出哪一侧整体不可能包含目标，或哪一侧整体已经满足可行性。二分的核心不是 mid 公式，而是被丢弃区间确实不含答案。

\textbf{状态复用。} 把逐个试候选改成维护答案所在区间；边界谓词、可行性检查或局部比较给出分支方向，每个分支都必须能排除一侧候选。本题落点是：规律是行尾小于下一行行首时可以整体展平二分，空矩阵和单行要先处理。手算时只改被新输入影响的那一小部分，而不是回头重算完整候选。

\textbf{指针和字段。} 状态字段要能说清：left、right、mid、mid 处比较值或检查返回值、答案区间含义、返回值语义；每一轮更新都要保留目标位置或第一个可行值。手算时按这个协议跟踪：拿矩阵 [[1, 3, 5], [7, 9, 11]] 找 9 手算，把它看成一维有序数组 [1, 3, 5, 7, 9, 11]，下标 4 映射到 row=1、col=1。然后按协议复查：手算时列出 left、mid、right，以及 mid 处比较结果、可行性检查结果或局部趋势。每一步都要写出被丢弃的一侧为什么不含答案。

\textbf{代码落点。} 检查函数或比较分支单独写清，mid 不溢出，循环退出条件和返回值配套；每个分支都要解释丢弃区间。关键代码行必须能逐句翻译成“旧状态怎样变成新状态”，否则就只是模板。

\textbf{正确性。} 证明从排除一侧开始：答案空间题证明谓词单调，局部比较题证明保留下来的区间仍含目标；不能只靠经验猜测左右边界。

\textbf{边界实验。} 先用空矩阵、单行、单列、目标小于全部或大于全部做反例检查。实验时覆盖空矩阵、单行、单列、目标小于全部或大于全部，逐轮打印 left、mid、right、比较值或检查结果。只要有一次被排除区间仍含答案，二分不变量就错了。

\textbf{复杂度变化。} 暴力逐个试候选是线性或和值域成正比；二分每次排除一半候选，复杂度变成对数级，答案二分还要乘上单次检查成本。

\textbf{迁移追问。} 如果题目改成“若只是行列分别有序但行之间不连续，应改用右上角排除法”，先判断原状态是否仍然覆盖新答案集合，再决定是微调边界、改 value 语义，还是换算法。

\noindent\textbf{LeetCode 81 搜索旋转排序数组 II。}

\textbf{题目说明。} LeetCode 81 搜索旋转排序数组 II 的题面入口要先还原成具体任务：旋转数组存在重复值时，二分可能无法判断哪一半有序。

\textbf{样例入口。} 拿 [2, 5, 6, 0, 0, 1, 2] 找 0 手算，可以判断一侧有序并排除；拿 [1, 1, 1, 1] 时 mid、left、right 相同，无法判断，只能收缩边界。

\textbf{暴力基线。} 慢版本按顺序扫描下标、逐个试答案值，或直接检查候选直到遇到目标边界。它先保证候选空间覆盖完整，但每次只排除一个候选，浪费了题目给出的顺序、比较或可行性边界。放到本题就是：先按“旋转数组存在重复值时，二分可能无法判断哪一半有序”完整试慢版本；样例里已经能看到“规律是重复元素让旋转二分可能退化，但每次丢掉一个相同边界不会漏目标”，所以优化前必须先能说清慢版本枚举了哪些候选。

\textbf{暴力过程。} 拿 [2, 5, 6, 0, 0, 1, 2] 找 0 手算，可以判断一侧有序并排除；拿 [1, 1, 1, 1] 时 mid、left、right 相同，无法判断，只能收缩边界。

\textbf{重复工作。} 题面模型：旋转数组存在重复值时，二分可能无法判断哪一半有序。暴力版的慢点要落到本题：慢版本会按顺序试“旋转数组存在重复值时，二分可能无法判断哪一半有序”里的候选；而样例规律已经指出“规律是重复元素让旋转二分可能退化，但每次丢掉一个相同边界不会漏目标”。一次比较或一次可行性检查能丢掉一整段候选，继续线性试就是浪费。后面的优化只消掉这类重复工作，不能改变答案集合。

\textbf{优化条件。} 本题的关键观察是：当左右和中点值相同导致信息不足时，收缩边界；否则仍按有序半边排除。这句话必须能翻译成可维护状态；如果题目条件改变后观察不再成立，就不能继续套原来的写法。

\textbf{相邻状态。} 规律是重复元素让旋转二分可能退化，但每次丢掉一个相同边界不会漏目标。把这个特例继续拆开看：每次比较 mid 之后，必须能指出哪一侧整体不可能包含目标，或哪一侧整体已经满足可行性。二分的核心不是 mid 公式，而是被丢弃区间确实不含答案。

\textbf{状态复用。} 把逐个试候选改成维护答案所在区间；边界谓词、可行性检查或局部比较给出分支方向，每个分支都必须能排除一侧候选。本题落点是：规律是重复元素让旋转二分可能退化，但每次丢掉一个相同边界不会漏目标。手算时只改被新输入影响的那一小部分，而不是回头重算完整候选。

\textbf{指针和字段。} 状态字段要能说清：left、right、mid、mid 处比较值或检查返回值、答案区间含义、返回值语义；每一轮更新都要保留目标位置或第一个可行值。手算时按这个协议跟踪：拿 [2, 5, 6, 0, 0, 1, 2] 找 0 手算，可以判断一侧有序并排除；拿 [1, 1, 1, 1] 时 mid、left、right 相同，无法判断，只能收缩边界。然后按协议复查：手算时列出 left、mid、right，以及 mid 处比较结果、可行性检查结果或局部趋势。每一步都要写出被丢弃的一侧为什么不含答案。

\textbf{代码落点。} 检查函数或比较分支单独写清，mid 不溢出，循环退出条件和返回值配套；每个分支都要解释丢弃区间。关键代码行必须能逐句翻译成“旧状态怎样变成新状态”，否则就只是模板。

\textbf{正确性。} 证明从排除一侧开始：答案空间题证明谓词单调，局部比较题证明保留下来的区间仍含目标；不能只靠经验猜测左右边界。

\textbf{边界实验。} 先用全相同元素、目标不存在、重复值跨过旋转点、退化到线性做反例检查。实验时覆盖全相同元素、目标不存在、重复值跨过旋转点、退化到线性，逐轮打印 left、mid、right、比较值或检查结果。只要有一次被排除区间仍含答案，二分不变量就错了。

\textbf{复杂度变化。} 暴力逐个试候选是线性或和值域成正比；二分每次排除一半候选，复杂度变成对数级，答案二分还要乘上单次检查成本。

\textbf{迁移追问。} 如果题目改成“这题说明二分依赖信息量，重复值会削弱每次排除一半的能力”，先判断原状态是否仍然覆盖新答案集合，再决定是微调边界、改 value 语义，还是换算法。

\noindent\textbf{LeetCode 153 寻找旋转排序数组中的最小值。}

\textbf{题目说明。} LeetCode 153 寻找旋转排序数组中的最小值 的题面入口要先还原成具体任务：最小值是旋转断点，右端点可帮助判断中点在哪一段。

\textbf{样例入口。} 拿 [3, 4, 5, 1, 2] 手算，mid=5 大于右端 2，说明最小值一定在 mid 右侧；若 mid 小于右端，最小值在左半含 mid。

\textbf{暴力基线。} 慢版本按顺序扫描下标、逐个试答案值，或直接检查候选直到遇到目标边界。它先保证候选空间覆盖完整，但每次只排除一个候选，浪费了题目给出的顺序、比较或可行性边界。放到本题就是：先按“最小值是旋转断点，右端点可帮助判断中点在哪一段”完整试慢版本；样例里已经能看到“规律是和右端比较能判断旋转点所在侧，未旋转数组会一路收缩到第一个元素”，所以优化前必须先能说清慢版本枚举了哪些候选。

\textbf{暴力过程。} 拿 [3, 4, 5, 1, 2] 手算，mid=5 大于右端 2，说明最小值一定在 mid 右侧；若 mid 小于右端，最小值在左半含 mid。

\textbf{重复工作。} 题面模型：最小值是旋转断点，右端点可帮助判断中点在哪一段。暴力版的慢点要落到本题：慢版本会按顺序试“最小值是旋转断点，右端点可帮助判断中点在哪一段”里的候选；而样例规律已经指出“规律是和右端比较能判断旋转点所在侧，未旋转数组会一路收缩到第一个元素”。一次比较或一次可行性检查能丢掉一整段候选，继续线性试就是浪费。后面的优化只消掉这类重复工作，不能改变答案集合。

\textbf{优化条件。} 本题的关键观察是：若 mid 大于 right，最小值在右侧；否则在左侧含 mid。这句话必须能翻译成可维护状态；如果题目条件改变后观察不再成立，就不能继续套原来的写法。

\textbf{相邻状态。} 规律是和右端比较能判断旋转点所在侧，未旋转数组会一路收缩到第一个元素。把这个特例继续拆开看：每次比较 mid 之后，必须能指出哪一侧整体不可能包含目标，或哪一侧整体已经满足可行性。二分的核心不是 mid 公式，而是被丢弃区间确实不含答案。

\textbf{状态复用。} 把逐个试候选改成维护答案所在区间；边界谓词、可行性检查或局部比较给出分支方向，每个分支都必须能排除一侧候选。本题落点是：规律是和右端比较能判断旋转点所在侧，未旋转数组会一路收缩到第一个元素。手算时只改被新输入影响的那一小部分，而不是回头重算完整候选。

\textbf{指针和字段。} 状态字段要能说清：left、right、mid、mid 处比较值或检查返回值、答案区间含义、返回值语义；每一轮更新都要保留目标位置或第一个可行值。手算时按这个协议跟踪：拿 [3, 4, 5, 1, 2] 手算，mid=5 大于右端 2，说明最小值一定在 mid 右侧；若 mid 小于右端，最小值在左半含 mid。然后按协议复查：手算时列出 left、mid、right，以及 mid 处比较结果、可行性检查结果或局部趋势。每一步都要写出被丢弃的一侧为什么不含答案。

\textbf{代码落点。} 检查函数或比较分支单独写清，mid 不溢出，循环退出条件和返回值配套；每个分支都要解释丢弃区间。关键代码行必须能逐句翻译成“旧状态怎样变成新状态”，否则就只是模板。

\textbf{正确性。} 证明从排除一侧开始：答案空间题证明谓词单调，局部比较题证明保留下来的区间仍含目标；不能只靠经验猜测左右边界。

\textbf{边界实验。} 先用未旋转、长度一、旋转一位、无重复条件做反例检查。实验时覆盖未旋转、长度一、旋转一位、无重复条件，逐轮打印 left、mid、right、比较值或检查结果。只要有一次被排除区间仍含答案，二分不变量就错了。

\textbf{复杂度变化。} 暴力逐个试候选是线性或和值域成正比；二分每次排除一半候选，复杂度变成对数级，答案二分还要乘上单次检查成本。

\textbf{迁移追问。} 如果题目改成“有重复时遇到相等只能收缩 right，最坏退化”，先判断原状态是否仍然覆盖新答案集合，再决定是微调边界、改 value 语义，还是换算法。

\noindent\textbf{LeetCode 154 寻找旋转排序数组中的最小值 II。}

\textbf{题目说明。} LeetCode 154 寻找旋转排序数组中的最小值 II 的题面入口要先还原成具体任务：重复元素会让 mid 和 right 相等时无法判断最小值方向。

\textbf{样例入口。} 拿 [2, 2, 2, 0, 1] 手算，mid 和 right 都是 2 时无法判断最小值在哪侧，只能 right-- 丢掉一个重复值。

\textbf{暴力基线。} 慢版本按顺序扫描下标、逐个试答案值，或直接检查候选直到遇到目标边界。它先保证候选空间覆盖完整，但每次只排除一个候选，浪费了题目给出的顺序、比较或可行性边界。放到本题就是：先按“重复元素会让 mid 和 right 相等时无法判断最小值方向”完整试慢版本；样例里已经能看到“规律是重复元素会削弱二分信息，最坏会退化成线性，但不会丢失最小值”，所以优化前必须先能说清慢版本枚举了哪些候选。

\textbf{暴力过程。} 拿 [2, 2, 2, 0, 1] 手算，mid 和 right 都是 2 时无法判断最小值在哪侧，只能 right-- 丢掉一个重复值。

\textbf{重复工作。} 题面模型：重复元素会让 mid 和 right 相等时无法判断最小值方向。暴力版的慢点要落到本题：慢版本会按顺序试“重复元素会让 mid 和 right 相等时无法判断最小值方向”里的候选；而样例规律已经指出“规律是重复元素会削弱二分信息，最坏会退化成线性，但不会丢失最小值”。一次比较或一次可行性检查能丢掉一整段候选，继续线性试就是浪费。后面的优化只消掉这类重复工作，不能改变答案集合。

\textbf{优化条件。} 本题的关键观察是：相等时只能安全地丢掉一个 right，因为它不是唯一必要候选。这句话必须能翻译成可维护状态；如果题目条件改变后观察不再成立，就不能继续套原来的写法。

\textbf{相邻状态。} 规律是重复元素会削弱二分信息，最坏会退化成线性，但不会丢失最小值。把这个特例继续拆开看：每次比较 mid 之后，必须能指出哪一侧整体不可能包含目标，或哪一侧整体已经满足可行性。二分的核心不是 mid 公式，而是被丢弃区间确实不含答案。

\textbf{状态复用。} 把逐个试候选改成维护答案所在区间；边界谓词、可行性检查或局部比较给出分支方向，每个分支都必须能排除一侧候选。本题落点是：规律是重复元素会削弱二分信息，最坏会退化成线性，但不会丢失最小值。手算时只改被新输入影响的那一小部分，而不是回头重算完整候选。

\textbf{指针和字段。} 状态字段要能说清：left、right、mid、mid 处比较值或检查返回值、答案区间含义、返回值语义；每一轮更新都要保留目标位置或第一个可行值。手算时按这个协议跟踪：拿 [2, 2, 2, 0, 1] 手算，mid 和 right 都是 2 时无法判断最小值在哪侧，只能 right-- 丢掉一个重复值。然后按协议复查：手算时列出 left、mid、right，以及 mid 处比较结果、可行性检查结果或局部趋势。每一步都要写出被丢弃的一侧为什么不含答案。

\textbf{代码落点。} 检查函数或比较分支单独写清，mid 不溢出，循环退出条件和返回值配套；每个分支都要解释丢弃区间。关键代码行必须能逐句翻译成“旧状态怎样变成新状态”，否则就只是模板。

\textbf{正确性。} 证明从排除一侧开始：答案空间题证明谓词单调，局部比较题证明保留下来的区间仍含目标；不能只靠经验猜测左右边界。

\textbf{边界实验。} 先用全相等、重复值跨断点、最小值出现多次做反例检查。实验时覆盖全相等、重复值跨断点、最小值出现多次，逐轮打印 left、mid、right、比较值或检查结果。只要有一次被排除区间仍含答案，二分不变量就错了。

\textbf{复杂度变化。} 暴力逐个试候选是线性或和值域成正比；二分每次排除一半候选，复杂度变成对数级，答案二分还要乘上单次检查成本。

\textbf{迁移追问。} 如果题目改成“这题说明二分依赖的信息不足时会退化，不是所有旋转题都严格对数”，先判断原状态是否仍然覆盖新答案集合，再决定是微调边界、改 value 语义，还是换算法。

\noindent\textbf{LeetCode 162 寻找峰值。}

\textbf{题目说明。} LeetCode 162 寻找峰值 的题面入口要先还原成具体任务：比较 mid 和 mid 加一 的斜率能判断某侧一定存在峰值。

\textbf{样例入口。} 拿 [1, 2, 3, 1] 手算，mid=2 时 nums[mid] 小于右邻 3，说明右侧上坡方向必有峰值；若大于右邻，左侧含 mid 必有峰值。

\textbf{暴力基线。} 慢版本按顺序扫描下标、逐个试答案值，或直接检查候选直到遇到目标边界。它先保证候选空间覆盖完整，但每次只排除一个候选，浪费了题目给出的顺序、比较或可行性边界。放到本题就是：先按“比较 mid 和 mid 加一 的斜率能判断某侧一定存在峰值”完整试慢版本；样例里已经能看到“规律是比较 mid 和 mid+1 就能保留一侧峰值区间，边界外可视为负无穷”，所以优化前必须先能说清慢版本枚举了哪些候选。

\textbf{暴力过程。} 拿 [1, 2, 3, 1] 手算，mid=2 时 nums[mid] 小于右邻 3，说明右侧上坡方向必有峰值；若大于右邻，左侧含 mid 必有峰值。

\textbf{重复工作。} 题面模型：比较 mid 和 mid 加一 的斜率能判断某侧一定存在峰值。暴力版的慢点要落到本题：慢版本会按顺序试“比较 mid 和 mid 加一 的斜率能判断某侧一定存在峰值”里的候选；而样例规律已经指出“规律是比较 mid 和 mid+1 就能保留一侧峰值区间，边界外可视为负无穷”。一次比较或一次可行性检查能丢掉一整段候选，继续线性试就是浪费。后面的优化只消掉这类重复工作，不能改变答案集合。

\textbf{优化条件。} 本题的关键观察是：上升则右侧有峰，下降则左侧含 mid 有峰。这句话必须能翻译成可维护状态；如果题目条件改变后观察不再成立，就不能继续套原来的写法。

\textbf{相邻状态。} 规律是比较 mid 和 mid+1 就能保留一侧峰值区间，边界外可视为负无穷。把这个特例继续拆开看：每次比较 mid 之后，必须能指出哪一侧整体不可能包含目标，或哪一侧整体已经满足可行性。二分的核心不是 mid 公式，而是被丢弃区间确实不含答案。

\textbf{状态复用。} 把逐个试候选改成维护答案所在区间；边界谓词、可行性检查或局部比较给出分支方向，每个分支都必须能排除一侧候选。本题落点是：规律是比较 mid 和 mid+1 就能保留一侧峰值区间，边界外可视为负无穷。手算时只改被新输入影响的那一小部分，而不是回头重算完整候选。

\textbf{指针和字段。} 状态字段要能说清：left、right、mid、mid 处比较值或检查返回值、答案区间含义、返回值语义；每一轮更新都要保留目标位置或第一个可行值。手算时按这个协议跟踪：拿 [1, 2, 3, 1] 手算，mid=2 时 nums[mid] 小于右邻 3，说明右侧上坡方向必有峰值；若大于右邻，左侧含 mid 必有峰值。然后按协议复查：手算时列出 left、mid、right，以及 mid 处比较结果、可行性检查结果或局部趋势。每一步都要写出被丢弃的一侧为什么不含答案。

\textbf{代码落点。} 检查函数或比较分支单独写清，mid 不溢出，循环退出条件和返回值配套；每个分支都要解释丢弃区间。关键代码行必须能逐句翻译成“旧状态怎样变成新状态”，否则就只是模板。

\textbf{正确性。} 证明从排除一侧开始：答案空间题证明谓词单调，局部比较题证明保留下来的区间仍含目标；不能只靠经验猜测左右边界。

\textbf{边界实验。} 先用长度一、峰在边界、多个峰、相邻不相等假设做反例检查。实验时覆盖长度一、峰在边界、多个峰、相邻不相等假设，逐轮打印 left、mid、right、比较值或检查结果。只要有一次被排除区间仍含答案，二分不变量就错了。

\textbf{复杂度变化。} 暴力逐个试候选是线性或和值域成正比；二分每次排除一半候选，复杂度变成对数级，答案二分还要乘上单次检查成本。

\textbf{迁移追问。} 如果题目改成“二维峰值需要按行列最大值进行更复杂二分”，先判断原状态是否仍然覆盖新答案集合，再决定是微调边界、改 value 语义，还是换算法。

\noindent\textbf{LeetCode 240 搜索二维矩阵 II。}

\textbf{题目说明。} LeetCode 240 搜索二维矩阵 II 的题面入口要先还原成具体任务：右上角同时具有向左变小、向下变大的排除能力。

\textbf{样例入口。} 拿矩阵 [[1, 4, 7], [2, 5, 8], [3, 6, 9]] 找 6 手算，从右上角 7 开始，7 太大就左移到 4，4 太小就下移到 5，再下移到 6。

\textbf{暴力基线。} 慢版本按顺序扫描下标、逐个试答案值，或直接检查候选直到遇到目标边界。它先保证候选空间覆盖完整，但每次只排除一个候选，浪费了题目给出的顺序、比较或可行性边界。放到本题就是：先按“右上角同时具有向左变小、向下变大的排除能力”完整试慢版本；样例里已经能看到“规律是右上角一列向下增大、一行向左减小，每步能排除一行或一列”，所以优化前必须先能说清慢版本枚举了哪些候选。

\textbf{暴力过程。} 拿矩阵 [[1, 4, 7], [2, 5, 8], [3, 6, 9]] 找 6 手算，从右上角 7 开始，7 太大就左移到 4，4 太小就下移到 5，再下移到 6。

\textbf{重复工作。} 题面模型：右上角同时具有向左变小、向下变大的排除能力。暴力版的慢点要落到本题：慢版本会按顺序试“右上角同时具有向左变小、向下变大的排除能力”里的候选；而样例规律已经指出“规律是右上角一列向下增大、一行向左减小，每步能排除一行或一列”。一次比较或一次可行性检查能丢掉一整段候选，继续线性试就是浪费。后面的优化只消掉这类重复工作，不能改变答案集合。

\textbf{优化条件。} 本题的关键观察是：当前值大于目标就左移，小于目标就下移。这句话必须能翻译成可维护状态；如果题目条件改变后观察不再成立，就不能继续套原来的写法。

\textbf{相邻状态。} 规律是右上角一列向下增大、一行向左减小，每步能排除一行或一列。把这个特例继续拆开看：每次比较 mid 之后，必须能指出哪一侧整体不可能包含目标，或哪一侧整体已经满足可行性。二分的核心不是 mid 公式，而是被丢弃区间确实不含答案。

\textbf{状态复用。} 把逐个试候选改成维护答案所在区间；边界谓词、可行性检查或局部比较给出分支方向，每个分支都必须能排除一侧候选。本题落点是：规律是右上角一列向下增大、一行向左减小，每步能排除一行或一列。手算时只改被新输入影响的那一小部分，而不是回头重算完整候选。

\textbf{指针和字段。} 状态字段要能说清：left、right、mid、mid 处比较值或检查返回值、答案区间含义、返回值语义；每一轮更新都要保留目标位置或第一个可行值。手算时按这个协议跟踪：拿矩阵 [[1, 4, 7], [2, 5, 8], [3, 6, 9]] 找 6 手算，从右上角 7 开始，7 太大就左移到 4，4 太小就下移到 5，再下移到 6。然后按协议复查：手算时列出 left、mid、right，以及 mid 处比较结果、可行性检查结果或局部趋势。每一步都要写出被丢弃的一侧为什么不含答案。

\textbf{代码落点。} 检查函数或比较分支单独写清，mid 不溢出，循环退出条件和返回值配套；每个分支都要解释丢弃区间。关键代码行必须能逐句翻译成“旧状态怎样变成新状态”，否则就只是模板。

\textbf{正确性。} 证明从排除一侧开始：答案空间题证明谓词单调，局部比较题证明保留下来的区间仍含目标；不能只靠经验猜测左右边界。

\textbf{边界实验。} 先用空矩阵、单行、单列、目标在角落、重复值做反例检查。实验时覆盖空矩阵、单行、单列、目标在角落、重复值，逐轮打印 left、mid、right、比较值或检查结果。只要有一次被排除区间仍含答案，二分不变量就错了。

\textbf{复杂度变化。} 暴力逐个试候选是线性或和值域成正比；二分每次排除一半候选，复杂度变成对数级，答案二分还要乘上单次检查成本。

\textbf{迁移追问。} 如果题目改成“从左上角无法一次排除一行或一列”，先判断原状态是否仍然覆盖新答案集合，再决定是微调边界、改 value 语义，还是换算法。

\noindent\textbf{LeetCode 315 计算右侧小于当前元素的个数。}

\textbf{题目说明。} LeetCode 315 计算右侧小于当前元素的个数 的题面入口要先还原成具体任务：从右向左扫描时，已处理元素正好是当前位置右侧元素集合。

\textbf{样例入口。} 拿 nums=[5, 2, 6, 1] 从右往左手算，先加入 1；处理 6 时右侧小于 6 的已有元素有 1 个；处理 2 时右侧小于 2 的也只有 1；处理 5 时右侧小于 5 的是 2 和 1。

\textbf{暴力基线。} 慢版本按顺序扫描下标、逐个试答案值，或直接检查候选直到遇到目标边界。它先保证候选空间覆盖完整，但每次只排除一个候选，浪费了题目给出的顺序、比较或可行性边界。放到本题就是：先按“从右向左扫描时，已处理元素正好是当前位置右侧元素集合”完整试慢版本；样例里已经能看到“规律是右侧元素集合动态增长，离散化后用树状数组查询当前排名之前的频次”，所以优化前必须先能说清慢版本枚举了哪些候选。

\textbf{暴力过程。} 拿 nums=[5, 2, 6, 1] 从右往左手算，先加入 1；处理 6 时右侧小于 6 的已有元素有 1 个；处理 2 时右侧小于 2 的也只有 1；处理 5 时右侧小于 5 的是 2 和 1。

\textbf{重复工作。} 题面模型：从右向左扫描时，已处理元素正好是当前位置右侧元素集合。暴力版的慢点要落到本题：慢版本会按顺序试“从右向左扫描时，已处理元素正好是当前位置右侧元素集合”里的候选；而样例规律已经指出“规律是右侧元素集合动态增长，离散化后用树状数组查询当前排名之前的频次”。一次比较或一次可行性检查能丢掉一整段候选，继续线性试就是浪费。后面的优化只消掉这类重复工作，不能改变答案集合。

\textbf{优化条件。} 本题的关键观察是：先离散化数值，再用树状数组维护右侧元素排名频次，查询当前排名之前的频次和。这句话必须能翻译成可维护状态；如果题目条件改变后观察不再成立，就不能继续套原来的写法。

\textbf{相邻状态。} 规律是右侧元素集合动态增长，离散化后用树状数组查询当前排名之前的频次。把这个特例继续拆开看：每次比较 mid 之后，必须能指出哪一侧整体不可能包含目标，或哪一侧整体已经满足可行性。二分的核心不是 mid 公式，而是被丢弃区间确实不含答案。

\textbf{状态复用。} 把逐个试候选改成维护答案所在区间；边界谓词、可行性检查或局部比较给出分支方向，每个分支都必须能排除一侧候选。本题落点是：规律是右侧元素集合动态增长，离散化后用树状数组查询当前排名之前的频次。手算时只改被新输入影响的那一小部分，而不是回头重算完整候选。

\textbf{指针和字段。} 状态字段要能说清：left、right、mid、mid 处比较值或检查返回值、答案区间含义、返回值语义；每一轮更新都要保留目标位置或第一个可行值。手算时按这个协议跟踪：拿 nums=[5, 2, 6, 1] 从右往左手算，先加入 1；处理 6 时右侧小于 6 的已有元素有 1 个；处理 2 时右侧小于 2 的也只有 1；处理 5 时右侧小于 5 的是 2 和 1。然后按协议复查：手算时列出 left、mid、right，以及 mid 处比较结果、可行性检查结果或局部趋势。每一步都要写出被丢弃的一侧为什么不含答案。

\textbf{代码落点。} 检查函数或比较分支单独写清，mid 不溢出，循环退出条件和返回值配套；每个分支都要解释丢弃区间。关键代码行必须能逐句翻译成“旧状态怎样变成新状态”，否则就只是模板。

\textbf{正确性。} 证明从排除一侧开始：答案空间题证明谓词单调，局部比较题证明保留下来的区间仍含目标；不能只靠经验猜测左右边界。

\textbf{边界实验。} 先用重复值、负数、值域很大、严格小于不能包含相等值、结果顺序要反向写回做反例检查。实验时覆盖重复值、负数、值域很大、严格小于不能包含相等值、结果顺序要反向写回，逐轮打印 left、mid、right、比较值或检查结果。只要有一次被排除区间仍含答案，二分不变量就错了。

\textbf{复杂度变化。} 暴力逐个试候选是线性或和值域成正比；二分每次排除一半候选，复杂度变成对数级，答案二分还要乘上单次检查成本。

\textbf{迁移追问。} 如果题目改成“若改成右侧大于当前元素，只需查询排名之后的频次”，先判断原状态是否仍然覆盖新答案集合，再决定是微调边界、改 value 语义，还是换算法。

\noindent\textbf{LeetCode 410 分割数组的最大值。}

\textbf{题目说明。} LeetCode 410 分割数组的最大值 的题面入口要先还原成具体任务：连续分成 m 段后最大段和越大，越容易完成分割。

\textbf{样例入口。} 拿 [7, 2, 5, 10, 8] 分成 2 段手算，若最大段和限制为 18，可以分成 [7, 2, 5] 和 [10, 8]；限制为 17 就需要三段。

\textbf{暴力基线。} 暴力枚举把数组切成 k 段的所有切法，计算每种切法的最大段和，取最小值。

\textbf{暴力过程。} 以 nums=[7, 2, 5, 10, 8]、k=2 为例，限制最大段和为 18 时可以切成 [7, 2, 5] 和 [10, 8]；限制为 17 时至少要三段。

\textbf{重复工作。} 题面模型：连续分成 m 段后最大段和越大，越容易完成分割。暴力版的慢点要落到本题：浪费在直接枚举切点。若给定一个容量 limit，可以线性判断是否能用不超过 k 段装下；limit 越大越容易成功。后面的优化只消掉这类重复工作，不能改变答案集合。

\textbf{优化条件。} 本题的关键观察是：二分最大段和上限，检查函数贪心累计，超过上限就新开一段。这句话必须能翻译成可维护状态；如果题目条件改变后观察不再成立，就不能继续套原来的写法。

\textbf{相邻状态。} 规律是答案空间是最大段和，可行性随容量增大单调变真，所以二分最小可行容量。把这个特例继续拆开看：每次比较 mid 之后，必须能指出哪一侧整体不可能包含目标，或哪一侧整体已经满足可行性。二分的核心不是 mid 公式，而是被丢弃区间确实不含答案。

\textbf{状态复用。} 答案范围是 [max(nums), sum(nums)]。二分 limit，check(limit) 贪心从左到右累加，超过 limit 就开新段。手算时只改被新输入影响的那一小部分，而不是回头重算完整候选。

\textbf{指针和字段。} 状态字段要能说清：limit 是候选最大段和，pieces 是当前需要段数，current\_sum 是当前段累计和。手算时按这个协议跟踪：limit=18 时，7+2+5=14，再加 10 超过，于是开第二段 10+8=18，段数 2 可行。

\textbf{代码落点。} 关键是 check 只回答可行性，不求具体最优切点；若 pieces<=k，说明 limit 可行，右边界左移。关键代码行必须能逐句翻译成“旧状态怎样变成新状态”，否则就只是模板。

\textbf{正确性。} 固定 limit 时，贪心尽量延长当前段会让段数最少；若最少段数仍超过 k，则任何切法都不可行。可行性随 limit 增大单调变真。

\textbf{边界实验。} 先用m 为一、m 等于数组长度、单个元素超过猜测上限、总和溢出做反例检查。实验时覆盖m 为一、m 等于数组长度、单个元素超过猜测上限、总和溢出，逐轮打印 left、mid、right、比较值或检查结果。只要有一次被排除区间仍含答案，二分不变量就错了。

\textbf{复杂度变化。} 枚举切法指数级或 DP O(nk)；答案二分 O(n log(sum-max)) 时间，O(1) 额外空间。

\textbf{迁移追问。} 如果题目改成“也可用 DP 求精确最优，但答案二分更适合大输入”，先判断原状态是否仍然覆盖新答案集合，再决定是微调边界、改 value 语义，还是换算法。

\noindent\textbf{LeetCode 540 有序数组中的单一元素。}

\textbf{题目说明。} LeetCode 540 有序数组中的单一元素 的题面入口要先还原成具体任务：除一个元素外其余元素都成对出现，单一元素会改变配对下标的奇偶规律。

\textbf{样例入口。} 拿 [1, 1, 2, 3, 3] 手算，单一元素 2 之前成对元素从偶数下标开始，之后配对位置会错位。

\textbf{暴力基线。} 慢版本按顺序扫描下标、逐个试答案值，或直接检查候选直到遇到目标边界。它先保证候选空间覆盖完整，但每次只排除一个候选，浪费了题目给出的顺序、比较或可行性边界。放到本题就是：先按“除一个元素外其余元素都成对出现，单一元素会改变配对下标的奇偶规律”完整试慢版本；样例里已经能看到“规律是用 mid 的偶偶配对关系判断单一元素在哪边，保持答案区间长度为奇数”，所以优化前必须先能说清慢版本枚举了哪些候选。

\textbf{暴力过程。} 拿 [1, 1, 2, 3, 3] 手算，单一元素 2 之前成对元素从偶数下标开始，之后配对位置会错位。

\textbf{重复工作。} 题面模型：除一个元素外其余元素都成对出现，单一元素会改变配对下标的奇偶规律。暴力版的慢点要落到本题：慢版本会按顺序试“除一个元素外其余元素都成对出现，单一元素会改变配对下标的奇偶规律”里的候选；而样例规律已经指出“规律是用 mid 的偶偶配对关系判断单一元素在哪边，保持答案区间长度为奇数”。一次比较或一次可行性检查能丢掉一整段候选，继续线性试就是浪费。后面的优化只消掉这类重复工作，不能改变答案集合。

\textbf{优化条件。} 本题的关键观察是：把 mid 调整到偶数位，若 nums[mid] 等于 nums[mid+1]，单一元素在右侧，否则在左侧。这句话必须能翻译成可维护状态；如果题目条件改变后观察不再成立，就不能继续套原来的写法。

\textbf{相邻状态。} 规律是用 mid 的偶偶配对关系判断单一元素在哪边，保持答案区间长度为奇数。把这个特例继续拆开看：每次比较 mid 之后，必须能指出哪一侧整体不可能包含目标，或哪一侧整体已经满足可行性。二分的核心不是 mid 公式，而是被丢弃区间确实不含答案。

\textbf{状态复用。} 把逐个试候选改成维护答案所在区间；边界谓词、可行性检查或局部比较给出分支方向，每个分支都必须能排除一侧候选。本题落点是：规律是用 mid 的偶偶配对关系判断单一元素在哪边，保持答案区间长度为奇数。手算时只改被新输入影响的那一小部分，而不是回头重算完整候选。

\textbf{指针和字段。} 状态字段要能说清：left、right、mid、mid 处比较值或检查返回值、答案区间含义、返回值语义；每一轮更新都要保留目标位置或第一个可行值。手算时按这个协议跟踪：拿 [1, 1, 2, 3, 3] 手算，单一元素 2 之前成对元素从偶数下标开始，之后配对位置会错位。然后按协议复查：手算时列出 left、mid、right，以及 mid 处比较结果、可行性检查结果或局部趋势。每一步都要写出被丢弃的一侧为什么不含答案。

\textbf{代码落点。} 检查函数或比较分支单独写清，mid 不溢出，循环退出条件和返回值配套；每个分支都要解释丢弃区间。关键代码行必须能逐句翻译成“旧状态怎样变成新状态”，否则就只是模板。

\textbf{正确性。} 证明从排除一侧开始：答案空间题证明谓词单调，局部比较题证明保留下来的区间仍含目标；不能只靠经验猜测左右边界。

\textbf{边界实验。} 先用单一元素在开头或结尾、数组长度一、mid 加一越界做反例检查。实验时覆盖单一元素在开头或结尾、数组长度一、mid 加一越界，逐轮打印 left、mid、right、比较值或检查结果。只要有一次被排除区间仍含答案，二分不变量就错了。

\textbf{复杂度变化。} 暴力逐个试候选是线性或和值域成正比；二分每次排除一半候选，复杂度变成对数级，答案二分还要乘上单次检查成本。

\textbf{迁移追问。} 如果题目改成“异或也能求值，但二分利用有序性达到对数时间”，先判断原状态是否仍然覆盖新答案集合，再决定是微调边界、改 value 语义，还是换算法。

\noindent\textbf{LeetCode 704 二分查找。}

\textbf{题目说明。} LeetCode 704 二分查找 的题面入口要先还原成具体任务：基础题的重点是固定区间语义。

\textbf{样例入口。} 拿 [1, 3, 5] 找 3 手算，mid 正好命中返回；找 4 时，3 太小丢左侧，5 太大丢右侧，区间为空返回 -1。

\textbf{暴力基线。} 暴力从左到右扫描数组，找到 nums[i]==target 就返回 i，否则返回 -1。

\textbf{暴力过程。} 以 nums=[1, 3, 5] 找 4 为例，暴力会看 1、3、5；但数组有序，看到 mid=3 小于 4 后，左半都可排除。

\textbf{重复工作。} 题面模型：基础题的重点是固定区间语义。暴力版的慢点要落到本题：浪费在有序数组里一次只排除一个元素。比较 mid 后，可以一次丢掉不可能包含 target 的半边。后面的优化只消掉这类重复工作，不能改变答案集合。

\textbf{优化条件。} 本题的关键观察是：左闭右开写法能统一 lower bound 和存在性检查。这句话必须能翻译成可维护状态；如果题目条件改变后观察不再成立，就不能继续套原来的写法。

\textbf{相邻状态。} 规律是标准二分每次排除一半，边界写法必须和退出条件匹配。把这个特例继续拆开看：每次比较 mid 之后，必须能指出哪一侧整体不可能包含目标，或哪一侧整体已经满足可行性。二分的核心不是 mid 公式，而是被丢弃区间确实不含答案。

\textbf{状态复用。} 采用左闭右闭或左闭右开都可以，但必须固定语义。左闭右闭写法中 left<=right，mid 命中返回；nums[mid]<target 时 left=mid+1，否则 right=mid-1。手算时只改被新输入影响的那一小部分，而不是回头重算完整候选。

\textbf{指针和字段。} 状态字段要能说清：left/right 是当前可能含 target 的区间边界，mid 是中点，nums[mid] 与 target 的比较决定排除哪侧。手算时按这个协议跟踪：[1, 3, 5] 找 4：mid 指向 3，目标更大，丢掉 1 和 3；mid 指向 5，目标更小，区间变空返回 -1。

\textbf{代码落点。} 关键是 mid=left+(right-left)/2 防溢出；循环条件和边界更新必须配套，不能混用闭区间和开区间写法。关键代码行必须能逐句翻译成“旧状态怎样变成新状态”，否则就只是模板。

\textbf{正确性。} 每轮被丢弃的一侧所有元素都因有序性严格小于或大于 target，不可能包含答案；保留区间仍覆盖所有可能位置。

\textbf{边界实验。} 先用空数组、长度一、目标在边界、目标不存在做反例检查。实验时覆盖空数组、长度一、目标在边界、目标不存在，逐轮打印 left、mid、right、比较值或检查结果。只要有一次被排除区间仍含答案，二分不变量就错了。

\textbf{复杂度变化。} 暴力 O(n)，二分 O(log n) 时间、O(1) 空间。

\textbf{迁移追问。} 如果题目改成“所有复杂二分都应该能退回到这道题的边界不变量”，先判断原状态是否仍然覆盖新答案集合，再决定是微调边界、改 value 语义，还是换算法。

\noindent\textbf{LeetCode 875 爱吃香蕉的珂珂。}

\textbf{题目说明。} LeetCode 875 爱吃香蕉的珂珂 的题面入口要先还原成具体任务：速度越大，总耗时越小，答案具有单调可行性。

\textbf{样例入口。} 拿 piles=[3, 6, 7, 11]、h=8 手算，速度 4 时四堆分别需要 1、2、2、3 小时，总计 8 小时，刚好可行；速度 3 会超过 8 小时。

\textbf{暴力基线。} 暴力从速度 1 一直试到 max(piles)，每个速度都计算吃完所有香蕉需要多少小时。

\textbf{暴力过程。} 以 piles=[3, 6, 7, 11]、h=8 为例，速度 4 时耗时 ceil(3/4)+ceil(6/4)+ceil(7/4)+ceil(11/4)=8，可行；速度 3 超时。

\textbf{重复工作。} 题面模型：速度越大，总耗时越小，答案具有单调可行性。暴力版的慢点要落到本题：浪费在速度是有序答案空间。速度越大，总耗时只会不增，所以可行性从 false 到 true 有单调边界。后面的优化只消掉这类重复工作，不能改变答案集合。

\textbf{优化条件。} 本题的关键观察是：二分速度，检查函数用向上取整累加小时数。这句话必须能翻译成可维护状态；如果题目条件改变后观察不再成立，就不能继续套原来的写法。

\textbf{相邻状态。} 规律是速度越大耗时越少，可行性单调，二分最小可行速度。把这个特例继续拆开看：每次比较 mid 之后，必须能指出哪一侧整体不可能包含目标，或哪一侧整体已经满足可行性。二分的核心不是 mid 公式，而是被丢弃区间确实不含答案。

\textbf{状态复用。} 二分最小可行速度。check(speed) 扫描 piles，累计每堆需要的小时数，判断是否 <=h。手算时只改被新输入影响的那一小部分，而不是回头重算完整候选。

\textbf{指针和字段。} 状态字段要能说清：left/right 是速度范围，mid 是候选速度，hours 是该速度下总耗时。手算时按这个协议跟踪：若 mid=6 可行，说明更大的速度也都可行，答案可能更小，收缩右边界；若 mid=3 不可行，左边界右移。

\textbf{代码落点。} 关键是整除向上取整写成 (pile+speed-1)/speed，并用 long long 累计 hours 防止溢出。关键代码行必须能逐句翻译成“旧状态怎样变成新状态”，否则就只是模板。

\textbf{正确性。} speed 增大时每个 ceil(pile/speed) 都不增，总小时数不增；因此可行谓词单调，二分第一个 true 正是最小速度。

\textbf{边界实验。} 先用速度下界、堆很大、h 等于堆数、乘法加法溢出做反例检查。实验时覆盖速度下界、堆很大、h 等于堆数、乘法加法溢出，逐轮打印 left、mid、right、比较值或检查结果。只要有一次被排除区间仍含答案，二分不变量就错了。

\textbf{复杂度变化。} 线性试速度 O(maxPile*n)，二分 O(n log maxPile)，空间 O(1)。

\textbf{迁移追问。} 如果题目改成“船运包裹和分割数组都是同类最小可行上限”，先判断原状态是否仍然覆盖新答案集合，再决定是微调边界、改 value 语义，还是换算法。

\noindent\textbf{LeetCode 981 基于时间的键值存储。}

\textbf{题目说明。} LeetCode 981 基于时间的键值存储 的题面入口要先还原成具体任务：同一个 key 下的记录按时间戳递增，查询是不超过目标时间的最新值。

\textbf{样例入口。} 拿 key=foo 的记录 (1, bar)、(4, baz)，查询 t=3 手算，答案应是时间不超过 3 的最后一条，也就是 bar。查询 t=4 才是 baz。

\textbf{暴力基线。} 慢版本按顺序扫描下标、逐个试答案值，或直接检查候选直到遇到目标边界。它先保证候选空间覆盖完整，但每次只排除一个候选，浪费了题目给出的顺序、比较或可行性边界。放到本题就是：先按“同一个 key 下的记录按时间戳递增，查询是不超过目标时间的最新值”完整试慢版本；样例里已经能看到“规律是每个 key 内部时间递增，二分第一个大于查询时间的位置，再回退一格”，所以优化前必须先能说清慢版本枚举了哪些候选。

\textbf{暴力过程。} 拿 key=foo 的记录 (1, bar)、(4, baz)，查询 t=3 手算，答案应是时间不超过 3 的最后一条，也就是 bar。查询 t=4 才是 baz。

\textbf{重复工作。} 题面模型：同一个 key 下的记录按时间戳递增，查询是不超过目标时间的最新值。暴力版的慢点要落到本题：慢版本会按顺序试“同一个 key 下的记录按时间戳递增，查询是不超过目标时间的最新值”里的候选；而样例规律已经指出“规律是每个 key 内部时间递增，二分第一个大于查询时间的位置，再回退一格”。一次比较或一次可行性检查能丢掉一整段候选，继续线性试就是浪费。后面的优化只消掉这类重复工作，不能改变答案集合。

\textbf{优化条件。} 本题的关键观察是：哈希表定位 key 后，在该 key 的时间列表中二分第一个大于目标时间的位置并回退一格。这句话必须能翻译成可维护状态；如果题目条件改变后观察不再成立，就不能继续套原来的写法。

\textbf{相邻状态。} 规律是每个 key 内部时间递增，二分第一个大于查询时间的位置，再回退一格。把这个特例继续拆开看：每次比较 mid 之后，必须能指出哪一侧整体不可能包含目标，或哪一侧整体已经满足可行性。二分的核心不是 mid 公式，而是被丢弃区间确实不含答案。

\textbf{状态复用。} 把逐个试候选改成维护答案所在区间；边界谓词、可行性检查或局部比较给出分支方向，每个分支都必须能排除一侧候选。本题落点是：规律是每个 key 内部时间递增，二分第一个大于查询时间的位置，再回退一格。手算时只改被新输入影响的那一小部分，而不是回头重算完整候选。

\textbf{指针和字段。} 状态字段要能说清：left、right、mid、mid 处比较值或检查返回值、答案区间含义、返回值语义；每一轮更新都要保留目标位置或第一个可行值。手算时按这个协议跟踪：拿 key=foo 的记录 (1, bar)、(4, baz)，查询 t=3 手算，答案应是时间不超过 3 的最后一条，也就是 bar。查询 t=4 才是 baz。然后按协议复查：手算时列出 left、mid、right，以及 mid 处比较结果、可行性检查结果或局部趋势。每一步都要写出被丢弃的一侧为什么不含答案。

\textbf{代码落点。} 检查函数或比较分支单独写清，mid 不溢出，循环退出条件和返回值配套；每个分支都要解释丢弃区间。关键代码行必须能逐句翻译成“旧状态怎样变成新状态”，否则就只是模板。

\textbf{正确性。} 证明从排除一侧开始：答案空间题证明谓词单调，局部比较题证明保留下来的区间仍含目标；不能只靠经验猜测左右边界。

\textbf{边界实验。} 先用key 不存在、查询早于第一条记录、相同 key 多次 set、时间戳递增假设做反例检查。实验时覆盖key 不存在、查询早于第一条记录、相同 key 多次 set、时间戳递增假设，逐轮打印 left、mid、right、比较值或检查结果。只要有一次被排除区间仍含答案，二分不变量就错了。

\textbf{复杂度变化。} 暴力逐个试候选是线性或和值域成正比；二分每次排除一半候选，复杂度变成对数级，答案二分还要乘上单次检查成本。

\textbf{迁移追问。} 如果题目改成“如果时间戳无序到达，必须先排序或改成有序表维护每个 key 的记录”，先判断原状态是否仍然覆盖新答案集合，再决定是微调边界、改 value 语义，还是换算法。

\noindent\textbf{LeetCode 1011 在 D 天内送达包裹。}

\textbf{题目说明。} LeetCode 1011 在 D 天内送达包裹 的题面入口要先还原成具体任务：容量越大，需要天数越少，存在最小可行容量。

\textbf{样例入口。} 拿 weights=[1, 2, 3, 4, 5]、D=3 手算，容量 6 可以分成 [1, 2, 3]、[4]、[5] 三天；容量 5 需要更多天。

\textbf{暴力基线。} 暴力从运载能力 1 开始试到 sum(weights)，每个能力模拟装船天数，找到第一个能在 days 内完成的能力。

\textbf{暴力过程。} weights=[1, 2, 3, 4, 5]、days=3 时，capacity=6 可分成 [1, 2, 3]、[4]、[5] 三天；capacity=5 则需要更多天。

\textbf{重复工作。} 题面模型：容量越大，需要天数越少，存在最小可行容量。暴力版的慢点要落到本题：浪费在能力值有单调性却线性试。运载能力越大，需要天数只会不增。后面的优化只消掉这类重复工作，不能改变答案集合。

\textbf{优化条件。} 本题的关键观察是：下界是最大包裹，上界是总重量，检查函数按顺序贪心装船。这句话必须能翻译成可维护状态；如果题目条件改变后观察不再成立，就不能继续套原来的写法。

\textbf{相邻状态。} 规律是容量越大所需天数越少，可行性单调，二分最小容量。把这个特例继续拆开看：每次比较 mid 之后，必须能指出哪一侧整体不可能包含目标，或哪一侧整体已经满足可行性。二分的核心不是 mid 公式，而是被丢弃区间确实不含答案。

\textbf{状态复用。} 答案范围是 [max(weights), sum(weights)]。二分 capacity，check(capacity) 贪心从左到右装，超载就开新一天。手算时只改被新输入影响的那一小部分，而不是回头重算完整候选。

\textbf{指针和字段。} 状态字段要能说清：capacity 是候选船载重，used\_days 是模拟所需天数，current\_load 是当天已装重量。手算时按这个协议跟踪：capacity=6 时，1+2+3 正好一船；再装 4 会超载，于是开第二天；最后 5 第三天，满足 days=3。

\textbf{代码落点。} 关键是 left 从最大单件重量开始，不能从 1 开始；check 返回 used\_days<=days 时，说明能力可行，right=mid。关键代码行必须能逐句翻译成“旧状态怎样变成新状态”，否则就只是模板。

\textbf{正确性。} 固定 capacity 时，尽量把当天装满的贪心会让天数最少；可行性随 capacity 增大单调变真，二分第一个 true。

\textbf{边界实验。} 先用D 为一、D 等于包裹数、单个包裹超过猜测容量做反例检查。实验时覆盖D 为一、D 等于包裹数、单个包裹超过猜测容量，逐轮打印 left、mid、right、比较值或检查结果。只要有一次被排除区间仍含答案，二分不变量就错了。

\textbf{复杂度变化。} 线性试能力 O(sum*n)，二分 O(n log(sum-max)) 时间，O(1) 空间。

\textbf{迁移追问。} 如果题目改成“它和分割数组最大值都是连续分段最大和模型”，先判断原状态是否仍然覆盖新答案集合，再决定是微调边界、改 value 语义，还是换算法。

\noindent\textbf{LeetCode 1044 最长重复子串。}

\textbf{题目说明。} LeetCode 1044 最长重复子串 的题面入口要先还原成具体任务：重复子串长度具有可行性单调：存在长度 L 的重复，则更短长度也存在。

\textbf{样例入口。} 拿 banana 手算，重复子串 ana 出现两次，长度 3 可行；若长度 4 不可行，则更长也不可能。

\textbf{暴力基线。} 慢版本按顺序扫描下标、逐个试答案值，或直接检查候选直到遇到目标边界。它先保证候选空间覆盖完整，但每次只排除一个候选，浪费了题目给出的顺序、比较或可行性边界。放到本题就是：先按“重复子串长度具有可行性单调：存在长度 L 的重复，则更短长度也存在”完整试慢版本；样例里已经能看到“规律是二分重复子串长度，用滚动哈希检查同长度子串是否重复，哈希相等时还要防冲突”，所以优化前必须先能说清慢版本枚举了哪些候选。

\textbf{暴力过程。} 拿 banana 手算，重复子串 ana 出现两次，长度 3 可行；若长度 4 不可行，则更长也不可能。

\textbf{重复工作。} 题面模型：重复子串长度具有可行性单调：存在长度 L 的重复，则更短长度也存在。暴力版的慢点要落到本题：慢版本会按顺序试“重复子串长度具有可行性单调：存在长度 L 的重复，则更短长度也存在”里的候选；而样例规律已经指出“规律是二分重复子串长度，用滚动哈希检查同长度子串是否重复，哈希相等时还要防冲突”。一次比较或一次可行性检查能丢掉一整段候选，继续线性试就是浪费。后面的优化只消掉这类重复工作，不能改变答案集合。

\textbf{优化条件。} 本题的关键观察是：二分长度，用滚动哈希检查是否存在两个相同长度子串哈希相同，并在冲突时确认字符串。这句话必须能翻译成可维护状态；如果题目条件改变后观察不再成立，就不能继续套原来的写法。

\textbf{相邻状态。} 规律是二分重复子串长度，用滚动哈希检查同长度子串是否重复，哈希相等时还要防冲突。把这个特例继续拆开看：每次比较 mid 之后，必须能指出哪一侧整体不可能包含目标，或哪一侧整体已经满足可行性。二分的核心不是 mid 公式，而是被丢弃区间确实不含答案。

\textbf{状态复用。} 把逐个试候选改成维护答案所在区间；边界谓词、可行性检查或局部比较给出分支方向，每个分支都必须能排除一侧候选。本题落点是：规律是二分重复子串长度，用滚动哈希检查同长度子串是否重复，哈希相等时还要防冲突。手算时只改被新输入影响的那一小部分，而不是回头重算完整候选。

\textbf{指针和字段。} 状态字段要能说清：left、right、mid、mid 处比较值或检查返回值、答案区间含义、返回值语义；每一轮更新都要保留目标位置或第一个可行值。手算时按这个协议跟踪：拿 banana 手算，重复子串 ana 出现两次，长度 3 可行；若长度 4 不可行，则更长也不可能。然后按协议复查：手算时列出 left、mid、right，以及 mid 处比较结果、可行性检查结果或局部趋势。每一步都要写出被丢弃的一侧为什么不含答案。

\textbf{代码落点。} 检查函数或比较分支单独写清，mid 不溢出，循环退出条件和返回值配套；每个分支都要解释丢弃区间。关键代码行必须能逐句翻译成“旧状态怎样变成新状态”，否则就只是模板。

\textbf{正确性。} 证明从排除一侧开始：答案空间题证明谓词单调，局部比较题证明保留下来的区间仍含目标；不能只靠经验猜测左右边界。

\textbf{边界实验。} 先用全相同字符、无重复、哈希冲突、长度边界、模数选择做反例检查。实验时覆盖全相同字符、无重复、哈希冲突、长度边界、模数选择，逐轮打印 left、mid、right、比较值或检查结果。只要有一次被排除区间仍含答案，二分不变量就错了。

\textbf{复杂度变化。} 暴力逐个试候选是线性或和值域成正比；二分每次排除一半候选，复杂度变成对数级，答案二分还要乘上单次检查成本。

\textbf{迁移追问。} 如果题目改成“也可用后缀数组或后缀自动机，适合系统学习字符串算法”，先判断原状态是否仍然覆盖新答案集合，再决定是微调边界、改 value 语义，还是换算法。

\noindent\textbf{LeetCode 1482 制作 m 束花所需的最少天数。}

\textbf{题目说明。} LeetCode 1482 制作 m 束花所需的最少天数 的题面入口要先还原成具体任务：天数越大，可用花越多，能制作的相邻花束数不会减少。

\textbf{样例入口。} 拿 bloomDay=[1, 10, 3, 10, 2]、m=3、k=1 手算，day=2 时开放两朵只能做两束，day=3 时开放三朵能做三束。若 k=2，就必须数相邻开放段。

\textbf{暴力基线。} 慢版本按顺序扫描下标、逐个试答案值，或直接检查候选直到遇到目标边界。它先保证候选空间覆盖完整，但每次只排除一个候选，浪费了题目给出的顺序、比较或可行性边界。放到本题就是：先按“天数越大，可用花越多，能制作的相邻花束数不会减少”完整试慢版本；样例里已经能看到“规律是天数越大可用花越多，检查函数扫描连续开放花，满 k 朵就做一束”，所以优化前必须先能说清慢版本枚举了哪些候选。

\textbf{暴力过程。} 拿 bloomDay=[1, 10, 3, 10, 2]、m=3、k=1 手算，day=2 时开放两朵只能做两束，day=3 时开放三朵能做三束。若 k=2，就必须数相邻开放段。

\textbf{重复工作。} 题面模型：天数越大，可用花越多，能制作的相邻花束数不会减少。暴力版的慢点要落到本题：慢版本会按顺序试“天数越大，可用花越多，能制作的相邻花束数不会减少”里的候选；而样例规律已经指出“规律是天数越大可用花越多，检查函数扫描连续开放花，满 k 朵就做一束”。一次比较或一次可行性检查能丢掉一整段候选，继续线性试就是浪费。后面的优化只消掉这类重复工作，不能改变答案集合。

\textbf{优化条件。} 本题的关键观察是：二分最小天数，检查函数扫描连续已开放花，满 k 朵就做一束并清零。这句话必须能翻译成可维护状态；如果题目条件改变后观察不再成立，就不能继续套原来的写法。

\textbf{相邻状态。} 规律是天数越大可用花越多，检查函数扫描连续开放花，满 k 朵就做一束。把这个特例继续拆开看：每次比较 mid 之后，必须能指出哪一侧整体不可能包含目标，或哪一侧整体已经满足可行性。二分的核心不是 mid 公式，而是被丢弃区间确实不含答案。

\textbf{状态复用。} 把逐个试候选改成维护答案所在区间；边界谓词、可行性检查或局部比较给出分支方向，每个分支都必须能排除一侧候选。本题落点是：规律是天数越大可用花越多，检查函数扫描连续开放花，满 k 朵就做一束。手算时只改被新输入影响的那一小部分，而不是回头重算完整候选。

\textbf{指针和字段。} 状态字段要能说清：left、right、mid、mid 处比较值或检查返回值、答案区间含义、返回值语义；每一轮更新都要保留目标位置或第一个可行值。手算时按这个协议跟踪：拿 bloomDay=[1, 10, 3, 10, 2]、m=3、k=1 手算，day=2 时开放两朵只能做两束，day=3 时开放三朵能做三束。若 k=2，就必须数相邻开放段。然后按协议复查：手算时列出 left、mid、right，以及 mid 处比较结果、可行性检查结果或局部趋势。每一步都要写出被丢弃的一侧为什么不含答案。

\textbf{代码落点。} 检查函数或比较分支单独写清，mid 不溢出，循环退出条件和返回值配套；每个分支都要解释丢弃区间。关键代码行必须能逐句翻译成“旧状态怎样变成新状态”，否则就只是模板。

\textbf{正确性。} 证明从排除一侧开始：答案空间题证明谓词单调，局部比较题证明保留下来的区间仍含目标；不能只靠经验猜测左右边界。

\textbf{边界实验。} 先用m*k 超过花数、k 为一、所有花同日开放、相邻性、乘法溢出做反例检查。实验时覆盖m*k 超过花数、k 为一、所有花同日开放、相邻性、乘法溢出，逐轮打印 left、mid、right、比较值或检查结果。只要有一次被排除区间仍含答案，二分不变量就错了。

\textbf{复杂度变化。} 暴力逐个试候选是线性或和值域成正比；二分每次排除一半候选，复杂度变成对数级，答案二分还要乘上单次检查成本。

\textbf{迁移追问。} 如果题目改成“若花束不要求相邻，检查函数会变成计数开放花，不再扫连续段”，先判断原状态是否仍然覆盖新答案集合，再决定是微调边界、改 value 语义，还是换算法。

\subsection{实验复盘}

追问通常会要求你讲清答案边界、可行性谓词或局部排除规则。请准备说明 left、right 的语义，为什么 mid 被排除或保留，以及返回值是否还需要二次验证。

复盘本节时先从 LeetCode 4 寻找两个正序数组的中位数、LeetCode 33 搜索旋转排序数组、LeetCode 34 排序数组中查找首尾位置 开始：每道题都先手写一个能覆盖全部候选的暴力版，再指出优化结构具体保存了哪一段历史信息，最后用相邻案例比较为什么不能互换解法。
